\clearpage\pagestyle{matSimple}\phantomsection\label{sources}\pdfbookmark[0]{Sources}{bm-sources}
\matPageTitre{Sources et édition}{Références officielles · France}
{\fontsize{9.5}{13.5}\selectfont \textbf{1. Programme d’enseignement de l’école maternelle}\par Arrêté publié au BO n° 41 du 31 octobre 2024, applicable à la rentrée 2025 ; version consolidée intégrant le BO n° 19 du 7 mai 2026. Domaine « L’acquisition des premiers outils mathématiques » et partie « Se repérer dans l’espace ».\par\href{https://eduscol.education.gouv.fr/sites/default/files/document/programme-cycle-1-consolide-127565.pdf}{\textcolor{otGreenDark}{Consulter la source officielle}}\par\medskip
\textbf{2. Ressources d’accompagnement Éduscol}\par Ressources pour le cycle 1 : construction du nombre, résolution de problèmes, formes et grandeurs, motifs organisés.\par\href{https://eduscol.education.gouv.fr/}{\textcolor{otGreenDark}{Consulter la source officielle}}\par\medskip
\textbf{3. Choix faits dans ce cahier}\par Les quantités travaillées ne dépassent pas quatre ; les collections comparées globalement sont très différentes (au moins du simple au double), comme le demande le programme pour les enfants de moins de quatre ans. Le nombre zéro n’est pas introduit. Les doigts sont montrés en commençant par le pouce, mais toute autre façon de montrer un nombre avec ses doigts est juste. Les solides ne sont pas dessinés : ils se manipulent en classe. Les assemblages de formes comportent au plus quatre pièces. Chaque fiche commence par une activité avec de vrais objets, décrite dans le guide de l’adulte.\par}
